\documentclass[11pt]{article}
\usepackage[T1]{fontenc}
\usepackage[a4paper,margin=2.2cm]{geometry}
\usepackage{amsmath,amssymb}
\usepackage[per-mode=symbol]{siunitx}
\usepackage{enumitem}
\usepackage{xcolor}
\usepackage[hidelinks]{hyperref}
% Auto-generated by paper/make_paper_assets.py - do not edit by hand.

\newcommand{\YieldL}{1608}
\newcommand{\PRL}{0.819}
\newcommand{\CUFL}{18.35}
\newcommand{\POAL}{1963}
\newcommand{\TdayL}{39.5}
\newcommand{\TmaxL}{57.3}
\newcommand{\TriseL}{10.1}
\newcommand{\RHL}{74.0}
\newcommand{\DewL}{1{,}163}
\newcommand{\SRL}{0.965}
\newcommand{\DegBaseL}{0.30}
\newcommand{\DegTHL}{0.34}
\newcommand{\DegPIDL}{0.15}
\newcommand{\DegCorrL}{0.09}
\newcommand{\DegTotL}{0.88}
\newcommand{\TOWL}{3{,}055}
\newcommand{\RHencL}{63.5}
\newcommand{\LifeL}{36{,}233}
\newcommand{\CapL}{80.9}
\newcommand{\LCOEL}{2.83}
\newcommand{\LossIncidenceAngleL}{3.44}
\newcommand{\LossSpectralL}{-0.23}
\newcommand{\LossSoilingL}{3.67}
\newcommand{\LossDewCondensationL}{0.02}
\newcommand{\LossModuleTemperatureL}{6.12}
\newcommand{\LossDcWiringMismatchMotionLidL}{3.46}
\newcommand{\LossInverterAndClippingL}{2.03}
\newcommand{\LossAvailabilityL}{1.00}
\newcommand{\YieldF}{1634}
\newcommand{\PRF}{0.837}
\newcommand{\CUFF}{18.66}
\newcommand{\POAF}{1953}
\newcommand{\TdayF}{36.5}
\newcommand{\TmaxF}{50.0}
\newcommand{\TriseF}{7.3}
\newcommand{\RHF}{77.5}
\newcommand{\DewF}{1{,}092}
\newcommand{\SRF}{0.988}
\newcommand{\DegBaseF}{0.30}
\newcommand{\DegTHF}{0.34}
\newcommand{\DegPIDF}{0.15}
\newcommand{\DegCorrF}{0.05}
\newcommand{\DegTotF}{0.84}
\newcommand{\TOWF}{3{,}719}
\newcommand{\RHencF}{68.4}
\newcommand{\LifeF}{37{,}004}
\newcommand{\CapF}{81.7}
\newcommand{\LCOEF}{3.54}
\newcommand{\LossIncidenceAngleF}{3.47}
\newcommand{\LossSpectralF}{-0.24}
\newcommand{\LossSoilingF}{1.27}
\newcommand{\LossDewCondensationF}{0.04}
\newcommand{\LossModuleTemperatureF}{4.78}
\newcommand{\LossDcWiringMismatchMotionLidF}{4.62}
\newcommand{\LossInverterAndClippingF}{2.03}
\newcommand{\LossAvailabilityF}{1.50}
\newcommand{\YieldS}{1590}
\newcommand{\PRS}{0.814}
\newcommand{\CUFS}{18.15}
\newcommand{\POAS}{1953}
\newcommand{\TdayS}{36.2}
\newcommand{\TmaxS}{49.4}
\newcommand{\TriseS}{7.0}
\newcommand{\RHS}{79.7}
\newcommand{\DewS}{1{,}389}
\newcommand{\SRS}{0.969}
\newcommand{\DegBaseS}{0.30}
\newcommand{\DegTHS}{0.36}
\newcommand{\DegPIDS}{0.21}
\newcommand{\DegCorrS}{0.30}
\newcommand{\DegTotS}{1.16}
\newcommand{\TOWS}{4{,}288}
\newcommand{\RHencS}{70.4}
\newcommand{\LifeS}{34{,}677}
\newcommand{\CapS}{75.5}
\newcommand{\LCOES}{4.43}
\newcommand{\LossIncidenceAngleS}{3.47}
\newcommand{\LossSpectralS}{-0.25}
\newcommand{\LossSoilingS}{3.29}
\newcommand{\LossDewCondensationS}{0.05}
\newcommand{\LossModuleTemperatureS}{4.64}
\newcommand{\LossDcWiringMismatchMotionLidS}{4.91}
\newcommand{\LossInverterAndClippingS}{2.03}
\newcommand{\LossAvailabilityS}{2.00}
\newcommand{\YieldLM}{1608}
\newcommand{\PRLM}{0.819}
\newcommand{\CUFLM}{18.35}
\newcommand{\POALM}{1963}
\newcommand{\TdayLM}{39.5}
\newcommand{\TmaxLM}{57.3}
\newcommand{\TriseLM}{10.1}
\newcommand{\RHLM}{74.0}
\newcommand{\DewLM}{1{,}163}
\newcommand{\SRLM}{0.965}
\newcommand{\DegBaseLM}{0.30}
\newcommand{\DegTHLM}{0.34}
\newcommand{\DegPIDLM}{0.02}
\newcommand{\DegCorrLM}{0.04}
\newcommand{\DegTotLM}{0.70}
\newcommand{\TOWLM}{3{,}055}
\newcommand{\RHencLM}{63.5}
\newcommand{\LifeLM}{37{,}002}
\newcommand{\CapLM}{84.5}
\newcommand{\LCOELM}{2.90}
\newcommand{\LossIncidenceAngleLM}{3.44}
\newcommand{\LossSpectralLM}{-0.23}
\newcommand{\LossSoilingLM}{3.67}
\newcommand{\LossDewCondensationLM}{0.02}
\newcommand{\LossModuleTemperatureLM}{6.12}
\newcommand{\LossDcWiringMismatchMotionLidLM}{3.46}
\newcommand{\LossInverterAndClippingLM}{2.03}
\newcommand{\LossAvailabilityLM}{1.00}
\newcommand{\YieldFM}{1634}
\newcommand{\PRFM}{0.837}
\newcommand{\CUFFM}{18.66}
\newcommand{\POAFM}{1953}
\newcommand{\TdayFM}{36.5}
\newcommand{\TmaxFM}{50.0}
\newcommand{\TriseFM}{7.3}
\newcommand{\RHFM}{77.5}
\newcommand{\DewFM}{1{,}092}
\newcommand{\SRFM}{0.988}
\newcommand{\DegBaseFM}{0.30}
\newcommand{\DegTHFM}{0.34}
\newcommand{\DegPIDFM}{0.01}
\newcommand{\DegCorrFM}{0.03}
\newcommand{\DegTotFM}{0.68}
\newcommand{\TOWFM}{3{,}719}
\newcommand{\RHencFM}{68.4}
\newcommand{\LifeFM}{37{,}698}
\newcommand{\CapFM}{84.9}
\newcommand{\LCOEFM}{3.60}
\newcommand{\LossIncidenceAngleFM}{3.47}
\newcommand{\LossSpectralFM}{-0.24}
\newcommand{\LossSoilingFM}{1.27}
\newcommand{\LossDewCondensationFM}{0.04}
\newcommand{\LossModuleTemperatureFM}{4.78}
\newcommand{\LossDcWiringMismatchMotionLidFM}{4.62}
\newcommand{\LossInverterAndClippingFM}{2.03}
\newcommand{\LossAvailabilityFM}{1.50}
\newcommand{\YieldSM}{1590}
\newcommand{\PRSM}{0.814}
\newcommand{\CUFSM}{18.15}
\newcommand{\POASM}{1953}
\newcommand{\TdaySM}{36.2}
\newcommand{\TmaxSM}{49.4}
\newcommand{\TriseSM}{7.0}
\newcommand{\RHSM}{79.7}
\newcommand{\DewSM}{1{,}389}
\newcommand{\SRSM}{0.969}
\newcommand{\DegBaseSM}{0.30}
\newcommand{\DegTHSM}{0.36}
\newcommand{\DegPIDSM}{0.02}
\newcommand{\DegCorrSM}{0.15}
\newcommand{\DegTotSM}{0.83}
\newcommand{\TOWSM}{4{,}288}
\newcommand{\RHencSM}{70.4}
\newcommand{\LifeSM}{36{,}045}
\newcommand{\CapSM}{81.9}
\newcommand{\LCOESM}{4.43}
\newcommand{\LossIncidenceAngleSM}{3.47}
\newcommand{\LossSpectralSM}{-0.25}
\newcommand{\LossSoilingSM}{3.29}
\newcommand{\LossDewCondensationSM}{0.05}
\newcommand{\LossModuleTemperatureSM}{4.64}
\newcommand{\LossDcWiringMismatchMotionLidSM}{4.91}
\newcommand{\LossInverterAndClippingSM}{2.03}
\newcommand{\LossAvailabilitySM}{2.00}
\newcommand{\GainYieldF}{1.6}
\newcommand{\GainYieldS}{-1.1}
\newcommand{\GainLifeF}{2.1}
\newcommand{\GainLifeS}{-4.3}
\newcommand{\GainLifeSM}{3.9}
\newcommand{\dTdayF}{3.0}
\newcommand{\dTmaxF}{7.3}
\newcommand{\WaterSaved}{11{,}700}
\newcommand{\GHI}{1{,}889}
\newcommand{\Tair}{28.1}
\newcommand{\RHair}{74.0}
\newcommand{\Rain}{1{,}043}
\newcommand{\SensLifeMin}{32{,}498}
\newcommand{\SensLifeMax}{37{,}005}
\newcommand{\SensDegMin}{0.78}
\newcommand{\SensDegMax}{1.49}
\newcommand{\SensSaltSpan}{2{,}854}
\newcommand{\SensMoistSpan}{1{,}674}
\newcommand{\SensDegPerK}{0.167}
\newcommand{\SensTOWZero}{2{,}253}
\newcommand{\SensTOWMax}{5{,}547}
\newcommand{\SensYieldLoPct}{4.6}
\newcommand{\BreakEvenSalinity}{$<$0.25}
\newcommand{\NSeeds}{20}
\newcommand{\EnsYieldF}{1.71}
\newcommand{\EnsYieldSdF}{0.08}
\newcommand{\EnsLifeF}{2.17}
\newcommand{\EnsLifeSdF}{0.10}
\newcommand{\EnsDegF}{-0.039}
\newcommand{\EnsDegSdF}{0.003}
\newcommand{\EnsYieldS}{-1.08}
\newcommand{\EnsYieldSdS}{0.05}
\newcommand{\EnsLifeS}{-4.32}
\newcommand{\EnsLifeSdS}{0.06}
\newcommand{\EnsDegS}{+0.288}
\newcommand{\EnsDegSdS}{0.002}
\newcommand{\EnsYieldLandMin}{1{,}570}
\newcommand{\EnsYieldLandMax}{1{,}649}
\newcommand{\EnsFreshWins}{20}

% Auto-generated by paper/make_paper_assets.py (review additions) - do not edit by hand.

\newcommand{\NMC}{1{,}000}
\newcommand{\NParams}{28}
\newcommand{\MCYieldMedF}{1.2}
\newcommand{\MCYieldLoF}{$-$1.7}
\newcommand{\MCYieldHiF}{5.3}
\newcommand{\MCYieldPosF}{78}
\newcommand{\MCLifeMedF}{1.5}
\newcommand{\MCLifeLoF}{$-$3.6}
\newcommand{\MCLifeHiF}{6.3}
\newcommand{\MCLifePosF}{75}
\newcommand{\MCYieldMedS}{$-$1.5}
\newcommand{\MCYieldLoS}{$-$5.5}
\newcommand{\MCYieldHiS}{3.1}
\newcommand{\MCYieldPosS}{27}
\newcommand{\MCLifeMedS}{$-$5.1}
\newcommand{\MCLifeLoS}{$-$11.5}
\newcommand{\MCLifeHiS}{0.7}
\newcommand{\MCLifePosS}{4}
\newcommand{\MCYieldMedSM}{$-$1.5}
\newcommand{\MCYieldLoSM}{$-$5.5}
\newcommand{\MCYieldHiSM}{3.1}
\newcommand{\MCYieldPosSM}{27}
\newcommand{\MCLifeMedSM}{$-$1.6}
\newcommand{\MCLifeLoSM}{$-$7.0}
\newcommand{\MCLifeHiSM}{3.8}
\newcommand{\MCLifePosSM}{27}
\newcommand{\RCNameA}{Land dust deposition}
\newcommand{\RCValA}{0.45}
\newcommand{\RCNameB}{FPV $U_1$}
\newcommand{\RCValB}{0.43}
\newcommand{\RCNameC}{Saline salt deposition}
\newcommand{\RCValC}{$-$0.41}
\newcommand{\RCNameD}{FPV $U_0$}
\newcommand{\RCValD}{0.30}
\newcommand{\RCNameE}{Saline chloride deposition}
\newcommand{\RCValE}{$-$0.28}
\newcommand{\RCNameF}{Land $U_1$}
\newcommand{\RCValF}{$-$0.26}
\newcommand{\RCFNameA}{Land dust deposition}
\newcommand{\RCFValA}{0.59}
\newcommand{\RCFNameB}{FPV $U_1$}
\newcommand{\RCFValB}{0.49}
\newcommand{\RCFNameC}{FPV $U_0$}
\newcommand{\RCFValC}{0.35}
\newcommand{\PLcoeF}{99.5}
\newcommand{\PLcoeS}{100.0}
\newcommand{\OATBase}{$-$4.30}
\newcommand{\ClStarText}{none in 0--400}
\newcommand{\SaltStarYield}{35}
\newcommand{\DTdStar}{2.7}
\newcommand{\LifeSalineClZeroPct}{$-$0.9}
\newcommand{\LifeSalineSaltZeroPct}{$-$0.9}
\newcommand{\YieldSalineSaltZeroPct}{1.6}
\newcommand{\DegSalineClMax}{1.31}
\newcommand{\DegSalineClZero}{0.87}
\newcommand{\ClMax}{400}
\newcommand{\TOWFreshAtStar}{5{,}745}
\newcommand{\DewFreshMin}{1.64}
\newcommand{\DewFreshMax}{1.66}
\newcommand{\DewSalineMin}{$-$1.10}
\newcommand{\DewSalineMax}{$-$1.06}
\newcommand{\MitGGvsLand}{1.5}
\newcommand{\MitAllvsLand}{3.3}
\newcommand{\MitPIDRec}{2.2}
\newcommand{\MitSaltRec}{1.7}
\newcommand{\MitBothRec}{3.9}
\newcommand{\MitBothvsLand}{$-$0.5}
\newcommand{\CapexStarF}{34.3}
\newcommand{\CapexStarS}{29.1}
\newcommand{\LandPremF}{10.5}
\newcommand{\LandPremS}{23.6}
\newcommand{\WaterStar}{93}
\newcommand{\LCOEWaterZero}{3.54}
\newcommand{\LCOEWaterTen}{3.47}
\newcommand{\LCOEWaterTwenty}{3.39}
\newcommand{\LCOEWaterForty}{3.24}
\newcommand{\LCOEWaterEighty}{2.93}
\newcommand{\FinGapFMin}{0.58}
\newcommand{\FinGapFMax}{0.86}
\newcommand{\FinGapSMin}{1.36}
\newcommand{\FinGapSMax}{1.87}

% Auto-generated: Monte Carlo convergence
\newcommand{\TabConv}{%
$\Delta E_{25}$, fresh & 1.49 (0.09) & 1.42 & [$-$3.57, 6.33] & [$-$3.55, 6.34] \\
$\Delta E_{25}$, saline & $-$5.08 (0.11) & $-$5.12 & [$-$11.50, 0.71] & [$-$11.22, 0.45] \\
$\Delta E_{25}$, saline (M) & $-$1.56 (0.09) & $-$1.65 & [$-$6.96, 3.78] & [$-$7.00, 3.65] \\
$\Delta E_{1}$, fresh & 1.21 (0.07) & 1.11 & [$-$1.67, 5.25] & [$-$1.74, 5.25] \\}
\newcommand{\ConvMaxShift}{0.28}
\newcommand{\ConvMaxSE}{0.30}
\newcommand{\NMCConv}{5{,}000}

\newcommand{\TOWsym}{\mathit{TOW}}
\sloppy
\setlength{\parindent}{0pt}
\setlength{\parskip}{5pt}
\newcounter{cmt}
\newenvironment{comment}[1]{\refstepcounter{cmt}\par\medskip\noindent\textbf{Comment \thecmt{} (#1).}\itshape\ }{\par}
\newenvironment{response}{\par\noindent\textbf{Response.}\ }{\par}
\newcommand{\change}[1]{\par\noindent\textcolor{blue!50!black}{\textbf{Changes:} #1}\par}

\begin{document}
\begin{center}
{\Large\bfseries Response to Reviewers --- Third Round}\\[4pt]
Manuscript: ``Lifetime Performance and Degradation of Floating Photovoltaic Systems Under Humid and Saline Conditions''
\end{center}

I thank the reviewer for the third assessment. This round implements the six priority revisions and the other
suggestions. None of the reported results changes: the convergence analysis reuses exactly the \NMC{} samples reported
in the previous version as the first part of a \NMCConv-sample run.

\section*{Priority 1: Monte Carlo convergence (\S21)}
\begin{comment}{convergence}
Show that 1,000 samples are sufficient, e.g.\ with 250--5,000 samples.
\end{comment}
\begin{response}
Done. A \NMCConv-sample Monte Carlo was run with the same seed. Its first \NMC{} samples are identical to the reported
analysis, so the statistics can be compared on nested subsets of 250, 500, 1,000, 2,000 and 5,000 samples, with bootstrap
standard errors. From \NMC{} to \NMCConv{} samples, no median or 95\% bound of the key outputs moves by more than
\ConvMaxShift~percentage point. The bootstrap standard errors at \NMC{} samples are at most \ConvMaxSE~percentage point,
which is small compared with the interval widths of 6--12 points.
\end{response}
\change{Section~III-I (method); Section~IV-H, Table~XII and Fig.~13 (new); \texttt{uncertainty.convergence()} and a new
unit test. The Monte Carlo now runs in parallel and is bit-identical to the serial run.}

\section*{Priority 2: Assumptions versus empirical evidence (\S6, \S7, \S19)}
\begin{comment}{model-evidence table}
Classify components into established models, literature-informed parameterizations and new assumptions.
\end{comment}
\begin{response}
Added as Table~III, introduced at the start of Section~III. It has three categories:
\begin{itemize}[nosep]
\item[A] Established models used as published: Spencer, Haurwitz, Erbs, Hay--Davies, ASHRAE IAM, Kasten--Young, Gueymard,
the First Solar spectral factor, Magnus, Faiman, PVWatts, the Arrhenius/Peck form, the ISO~9223 concepts and LCOE.
\item[B] Literature-informed parameter values: heat-loss coefficients, Peck $E_a$ and $n$, PID activation energy, NaCl
deliquescence RH, ISO~9223 chloride classes, the dust coefficient, and reference rates calibrated to the magnitude of
field surveys.
\item[C] New modelling assumptions: the over-water microclimate, the surface-RH/encapsulant-lag formulation, the dew
criterion and loss, salt deposition and cementation, the salt-assisted leakage function, the corrosion scaling, additive
constant-rate degradation, the mitigation factors, and all deposition rates and costs.
\end{itemize}
The text notes that the contribution lies in the coupling and in category~C, which is also where field calibration is
most needed.
\end{response}

\begin{comment}{PID submodel}
State that the PID submodel is a physically motivated scenario model, not an experimentally calibrated predictive model.
\end{comment}
\begin{response}
The sentence has been added to Section~III-F in the reviewer's words. For readability, Eq.~(13) now defines the
shifted threshold $\mathit{RH}_{50}^{*}=\mathit{RH}_{50}-\Delta_s s$ and the width $w$ explicitly.
\end{response}

\begin{comment}{corrosion positioning}
Position the corrosion term carefully as a calibration-based representation.
\end{comment}
\begin{response}
Section~III-F now states that Eq.~(14) ``borrows the ISO~9223 variables but not a fitted damage function''. It adds that
the linear time-of-wetness and square-root chloride scaling make it ``a calibrated scenario representation of corrosion,
not a validated corrosion-damage model''.
\end{response}

\section*{Priority 3: Validation positioning (\S3, \S10, \S11, \S17)}
\begin{comment}{wording}
Use ``baseline modelled cooling''; call Table~III a ``synthetic-weather consistency check''; keep the scenario framing;
emphasize the validation protocol.
\end{comment}
\begin{response}
All are done:
\begin{itemize}[nosep]
\item Section~IV-A now reads ``The baseline modelled cooling is \dTdayF~K in daytime \ldots{} This is a consequence of
the assumed FPV heat-loss coefficients, not a measured cooling effect.'' The Discussion and conclusion use the same
term.
\item The weather table (now Table~IV) is titled ``Synthetic-Weather Consistency Check''.
\item The field-validation protocol is moved forward to Section~V-C, directly after the decision-relevant drivers it is
designed to measure. It is also listed as a contribution in Section~I.
\end{itemize}
\end{response}

\section*{Priority 4: Constant-rate lifetime extrapolation (\S12)}
\begin{comment}{constant rate}
Add a sentence immediately before Eq.~(15) on equivalent constant-rate projections.
\end{comment}
\begin{response}
Added immediately before Eq.~(15): ``Consequently, the lifetime estimates should be interpreted as equivalent
constant-rate projections rather than predictions of the exact annual degradation trajectory.''
\end{response}

\section*{Priority 5: Baseline versus Monte Carlo labelling (\S18)}
\begin{comment}{labelling}
Add the word ``baseline'' every time the $+2.1$\% and $-4.3$\% values appear.
\end{comment}
\begin{response}
Done. Every occurrence in the abstract, Sections~IV-B, IV-E and IV-H, and the conclusion is now labelled ``baseline''.
The Monte Carlo medians ($+$\MCLifeMedF\% and \MCLifeMedS\%) are always reported with their 95\% intervals.
\end{response}

\section*{Priority 6: Equation and typesetting audit (\S20)}
\begin{comment}{typesetting}
Check Eqs.~(9)--(14) and all subscripts, superscripts and units in the final PDF.
\end{comment}
\begin{response}
The compiled PDF was inspected at high resolution. Equations~(9)--(14) render correctly with IEEEtran, so the compression
the reviewer saw is most likely a PDF-to-text extraction artefact. Three equations were nevertheless restructured for
readability:
\begin{itemize}[nosep]
\item Eq.~(13) is split over two lines with a named shifted threshold.
\item Eq.~(15) is written with a single sum $\sum(O_y-B_y)(1+r)^{-y}$, so that its equation number stays on the same
line.
\item The ambient-RH condition in the definition of $\TOWsym$ is stated explicitly.
\end{itemize}
The final PDF compiles without overfull boxes, and $T_a'$, $T_m$ and the RH subscripts are used consistently.
\end{response}

\section*{Other comments}
\begin{comment}{\S14, threshold wording}
Present $\Delta T_d^*\approx2.7$~K as a model-predicted site-screening criterion, not a universal threshold.
\end{comment}
\begin{response}
Section~IV-F now says it ``is a threshold predicted by this model under the stated assumptions, not a universal physical
constant; it can nevertheless serve as a site-screening criterion''. The abstract, Section~V-B and the conclusion say
``in this model''.
\end{response}

\begin{comment}{\S22, regional economics}
Phrase economic conclusions as ``under the assumed Indian-market cost and financing conditions''.
\end{comment}
\begin{response}
Done in Section~IV-I (opening sentence and the justification sentence).
\end{response}

\begin{comment}{\S4, \S13, preserve wording}
Retain the ``model-based scenario predictions'' and ``scenario-based mitigation benefits'' language.
\end{comment}
\begin{response}
Retained unchanged, and not strengthened.
\end{response}

\section*{Remaining limitation}
Measured FPV and site-weather data are still unavailable, so the manuscript remains a physically informed scenario and
uncertainty framework. Its evidential basis is now classified explicitly (Table~III), and a field-validation protocol is
proposed (Table~XIV).

\end{document}
